\cleardoublepage
\chapter*{Nota Histórica --- Historical Note --- 历史注记}
\label{nota-historica}
\addcontentsline{toc}{chapter}{Nota Histórica · Historical Note · 历史注记}
\markboth{Nota Histórica · Historical Note · 历史注记}{}

\thispagestyle{empty}

\noindent\textbf{Português.} Este livro é uma obra de ficção. Joaquim Antônio Correia, o Paciente 1.260, e a criatura que ele carrega não existiram. A instituição narrada como \textit{o Colônia} é ficcional. Mas o solo sobre o qual esta ficção foi construída é real --- e é por isso que ela precisa terminar aqui, e não com um susto.

\vspace{0.6em}
\noindent Na seca de 1915, o governo do Ceará criou o \textbf{Campo do Alagadiço}, próximo a Fortaleza, onde retirantes foram cercados e confinados. O \textbf{Campo do Patu}, em Senador Pompeu, pertence a outro momento histórico: a seca de 1932, com funcionamento situado em 1932--1933. Esta edição adota Alagadiço como referente do cenário de 1915 e mantém Joaquim, sua família, a vala e a transferência de 1917 como elementos ficcionais. Os totais de internados e mortos não são apresentados como precisos.

\vspace{0.6em}
\noindent O \textbf{Hospital Colônia de Barbacena}, em Minas Gerais, existiu de verdade. Ao longo do século~XX, estima-se que cerca de \textbf{60~mil pessoas} morreram entre seus muros --- muitas delas sem diagnóstico psiquiátrico algum: eram pobres, negras, epilépticas, mães solteiras, homossexuais, ``inconvenientes''. Houve eletrochoque sem anestesia, contenção, fome, valas comuns e a venda de corpos a faculdades de medicina. Nada disso é invenção deste livro; é história documentada --- ver, entre outras fontes, \textit{Holocausto Brasileiro}, de Daniela Arbex.

\vspace{0.6em}
\noindent O horror sobrenatural destas páginas é uma metáfora. A dor que atravessa gerações, que ninguém processa e que passa adiante como um contágio, é a única criatura verdadeira do livro --- e ela tem nome: esquecimento. O pior que se pode fazer com os mortos dos currais e da Colônia não é assombrá-los com uma história. É deixá-los sem nenhuma. Aos leitores afetados: interromper a leitura é sempre legítimo, e o apoio está disponível --- CVV 188 (24 horas, Brasil) e 988 Suicide \& Crisis Lifeline (Estados Unidos).

\vspace{1.2em}
\noindent\rule{\textwidth}{0.5pt}
\vspace{0.8em}

\noindent\textbf{English.} This book is a work of fiction. Joaquim Antônio Correia, Patient 1,260, and the creature he carries never existed. The institution narrated as \textit{the Colony} is fictional. But the ground upon which this fiction was built is real --- and that is why it must end here, and not with a jump scare.

\vspace{0.6em}
\noindent During the 1915 drought, the government of Ceará created the \textbf{Alagadiço Camp}, near Fortaleza, where retirantes were fenced in and confined. The \textbf{Patu Camp}, in Senador Pompeu, belongs to a different historical moment: the 1932 drought, with operation situated in 1932--1933. This edition adopts Alagadiço as the referent for its 1915 setting and retains Joaquim, his family, the trench, and the 1917 transfer as fictional elements. Totals of internees and deaths are not presented as precise.

\vspace{0.6em}
\noindent The \textbf{Barbacena Hospital-Colony}, in Minas Gerais, truly existed. Over the course of the twentieth century, an estimated \textbf{60,000 people} died within its walls --- many of them with no psychiatric diagnosis whatsoever: they were poor, Black, epileptic, single mothers, homosexuals, ``inconvenient'' people. There was electroshock without anesthesia, restraint, hunger, mass graves, and the sale of bodies to medical schools. None of this is invented by this book; it is documented history --- see, among other sources, \textit{Brazilian Holocaust}, by Daniela Arbex.

\vspace{0.6em}
\noindent The supernatural horror of these pages is a metaphor. The pain that crosses generations, that no one processes and that passes on like a contagion, is the only true creature of the book --- and it has a name: forgetting. The worst thing one can do with the dead of the pens and the Colony is not to haunt them with a story. It is to leave them with none. To the reader who feels affected: interrupting the reading is always legitimate, and support is available --- CVV 188 (24 hours, Brazil) and the 988 Suicide \& Crisis Lifeline (United States).

\vspace{1.2em}
\noindent\rule{\textwidth}{0.5pt}
\vspace{0.8em}

\noindent\textbf{中文.} 本书是一部虚构作品。若阿金·安东尼奥·科雷亚、患者1,260和他携带的怪物从未存在。被叙述为\textit{收容院}的机构是虚构的。但这一切虚构所依托的地基是真实的——正因如此，它必须在这里结束，而不是以一次惊吓结束。

\vspace{0.6em}
\noindent 1915年大旱期间，塞阿拉政府在福塔莱萨附近建立了\textbf{阿拉加迪索（Alagadiço）收容营}，逃荒者在那里遭到围禁。塞纳多尔·蓬佩乌的\textbf{帕图（Patu）收容营}属于另一个历史阶段：1932年大旱，其运作时间为1932--1933年。本版以阿拉加迪索作为1915年场景的历史参照；若阿金、他的家人、乱葬坑以及1917年的转送均为虚构元素。收容与死亡人数不作为精确数字呈现。

\vspace{0.6em}
\noindent 米纳斯吉拉斯州的\textbf{巴巴塞纳收容院}真实存在。二十世纪期间，估计有\textbf{六万人}在其墙内死亡——其中许多人根本没有精神科诊断：他们是穷人、黑人、癫痫患者、单亲母亲、同性恋者、“碍事”的人。有无麻醉的电休克、束缚、饥饿、乱葬坑，以及向医学院出售遗体。这些都不是本书杜撰的；它们是有据可查的历史——参见达妮埃拉·阿尔贝克斯的\textit{巴西大屠杀}等资料。

\vspace{0.6em}
\noindent 这些书页的超自然恐怖是一种隐喻。那种跨越世代、无人处理、像传染病一样传递的痛苦，才是本书唯一的真实怪物——它有名字：遗忘。对牲畜圈和收容院死者最糟糕的做法，不是用一个故事去闹鬼。而是不给他们留下任何故事。对感到受影响的读者：中断阅读永远是正当的，支持渠道可用——CVV 188（巴西，24小时）、988 自杀与危机生命线（美国），以及中国心理援助热线 12356。

\vspace{1.2em}

\noindent\hfill\textit{Que o que o tempo sepultou não seja apagado. --- So that what time buried may not be erased. --- 愿时间埋葬的，不被抹去。}\hspace{2em}
